
\vspace{3\baselineskip}

\begin{center}
  {\large\bfseries Abstract}
\end{center}
\vspace{0.25\baselineskip}

\begin{adjustwidth}{0.05\textwidth}{0.05\textwidth} % -> 0.90\textwidth block
\noindent
Amidst the biodiversity crisis, there is high demand for spatially explicit biodiversity monitoring. Global models that quantify impacts of human pressures provide important insights for conservation, but their accuracy in spatial projections has yet to be systematically tested. Here we evaluate this using a global dataset of 25,987 species inventories from 681 studies. Despite estimated land-use impacts in line with previous research, our results highlight the challenging gap between effect size inference and prediction. We find that mixed models with study attributes as random effects -- common in meta-analysis and used in several indicators -- exhibit generally low predictive accuracy. This is driven by reliance on a small set of averaged fixed effects. In contrast, a biogeographic-taxonomic model structure with explicit environmental covariates shows higher but still modest interpolation accuracy. However, performance when extrapolating to other contexts remains low, due to distribution shifts in environmental conditions. These patterns apply to site-level diversity and differences between sites. Models are essential for informed conservation efforts, but their applicability is fundamentally constrained by data availability. Whereas countries with extensive data can build high-fidelity national indicators, accelerated data collection and model development are needed to better support data-poor regions with localized, actionable biodiversity insights.
\end{adjustwidth}
